%%
%% fall-lecture21.tex
%% 
%% Made by Alex Nelson <pqnelson@gmail.com>
%% Login   <alex@lisp>
%% 
%% Started on  2026-11-15T09:47:32-0800
%% Last update 2026-11-15T09:47:32-0800
%% 

\lecture{}


\begin{node}
We have discussed~(\S\ref{node:fall2026:lec20:math157a:construction-omega-0}) our construction of a bundle $\omega_{0}\to E_{0}$
for a complex rank-$(n-1)$ bundle over $E_{0}$ which uses the
Hermitian metric on $E$, but this construction can be generalized to
avoid using the Hermitian metric.
\end{node}

\begin{node}
We have also discussed the Gysin sequence, which is also discussed briefly in the textbook~\cite[\S12.1]{milnor1974characteristic}.
Recall this looked like:
\begin{equation}
\dots\to H^{j-2n}(B)\xrightarrow{\smile e(\omega_{\RR})}H^{j}(B)\xrightarrow{\pi_{0}^{*}}H^{j}(E_{0})\xrightarrow{\varphi^{-1}\circ\delta}H^{j-2n+1}(B)\to\dots
\end{equation}
where we suppress the coefficients (they're all usually $\ZZ$ or $\ZZ_{2}$),
and for $0\leq j\leq 2n-2$ we say $\pi_{0}^{*}$ is an isomorphism.
\end{node}

\begin{puzzle}
Can we use a similar construction process [similar to the Chern
classes] to obtain the Stiefel--Whitney classes?
\end{puzzle}

\begin{node}
We can make the attempt! Let $\xi$ be an $\RR^{n}$-bundle. We start
with the top Stiefel--Whitney class
\begin{equation}
w_{n}(\xi):=e_{2}(\xi)\in H^{n}(B;\ZZ_{2})
\end{equation}
where $e_{2}(\xi)$ is the mod-2 Euler class.
\end{node}

\begin{node}
Now, we remove the zero section of $E$ to get $E_{0}(\xi)$, and we
construct the bundle $\xi_{0}\xrightarrow{\pi_{0}} E_{0}$ over $E_{0}$ analogous to our
construction~(\S\ref{node:fall2026:lec20:math157a:construction-omega-0})
for $\omega_{0}$. This $\xi_{0}$ is a real rank-$(n-1)$ bundle.
\end{node}

\begin{node}
The analogous Gysin sequence will take $\ZZ_{2}$
coefficients\footnote{Of course, oriented real vector bundles could
extend this to cohomology groups with integer coefficients. See~\cite[Ch~17~Thm~7.4]{husemoller1993fibre}.} and
looks like
\begin{equation}
\dots\to H^{j-2n}(B)\xrightarrow{\smile e_{2}(\xi)}H^{j}(B)\xrightarrow{\pi_{0}^{*}}H^{j}(E_{0})\xrightarrow{\varphi^{-1}\circ\delta}H^{j-2n+1}(B)\to\dots
\end{equation}
and for $0\leq j\leq n-2$ we see that $\pi_{0}^{*}$ is an isomorphism.
(Note: for \emph{unoriented} bundles, there is a generalization of the
Gysin sequence using $\smile w_{n}(\xi)$ instead of $\smile e_{2}(\xi)$.)
So what about $j=n-1$?

At $j=n-1$, the Gysin sequence looks like:
\begin{equation}
0\to H^{n-1}(B;\ZZ_{2})\xrightarrow{\pi_{0}^{*}}H^{n-1}(E_{0};\ZZ_{2})\xrightarrow{\varphi^{-1}\circ\delta}H^{0}(B;\ZZ_{2})\to\dots.
\end{equation}
We need to show
\begin{equation}
(\varphi^{-1}\circ\delta)(w_{n-1}(\xi_{0}))\stackrel{?}{=}0,
\end{equation}
then we have a unique image with $\pi_{0}^{*}$ and then we can use the same
argument as before. It turns out the Gysin sequence is related to the
sphere bundle (which is germane to persuing this question).
\end{node}

\begin{node}
For complex vector bundles, $S(\omega)=\{v\in E(\omega)\mid\mu(v)=1\}$
is a bundle. In the real case, we have $E_{0}\to S(\xi)$ we can
replace $H^{n-1}(E_{0};\ZZ_{2})$ by $H^{n-1}(S(\xi);\ZZ_{2})$. We need
to use the sphere bundle version of Gysin sequences.
\end{node}

\begin{recall}
Our goal is still to prove $\varphi\circ\delta(w_{n-1}(\xi))=0$.
This is nontrivial.
\end{recall}

\begin{node}
Let us take a look at the fiber of $S(\xi)\to B$, which is
$S(F_{b})\iso\sphere{n-1}$ where $F_{b}=\pi^{-1}\{b\}$ is the fiber
over $b\in B$ of $\xi$. Then
\begin{equation}
\xi_{0}|_{S(F_{b})}\iso\T S(F_{b})\quad\mbox{homeomorphic},
\end{equation}
since $v^{\perp}\iso\T_{v}S(F_{b})$ is the tangent plane to the vector
$v$ in the sphere.
\end{node}

\begin{node}
Without loss of generality, $B$ is path-connected.
\end{node}

\begin{node}
Now, recall the fiber inclusion
\begin{equation}
j_{b}\colon(D(F_{b}),S(F_{b}))\to(D(\xi),S(\xi)),
\end{equation}
so
\begin{equation}
i_{b}\colon S(F_{b})\to S(F).
\end{equation}
We also have the [unique] mod-2 Thom class $U_{b}\in H^{n}(D(F_{b}),S(F_{b});\ZZ_{2})$
indexed by $b\in B$. Then we have the commutative diagram
\begin{equation}
\vcenter{\xymatrix{H^{n-1}(S(F_{b});\ZZ_{2})\ar[r]^-{\delta} & H^{n}(D(F_{b}),S(F_{b});\ZZ_{2})\\
H^{n-1}(S(\xi);\xi)\ar[u]^{(i_{b})^{*}}\ar[r]^-{\delta} & H^{n}(D(\xi),S(\xi);\ZZ_{2})\ar[u]_{(j_{b})^{*}}}}
\end{equation}
Since $S(\xi)$ deformation retracts to $E_{0}$, $w_{n-1}(\xi_{0})\in H^{n-1}(S(\xi);\ZZ_{2})$,
so
\begin{subequations}
  \begin{align}
(i_{b})^{*}w_{n-1}(\xi_{0}) &= w_{n-1}(\xi_{0}|_{S(F_{b})})\\
&= w_{n-1}(\T\sphere{n-1}).
  \end{align}
\end{subequations}
Then using Theorem~\ref{thm:fall2026:math157a:very-important:lec18},
\begin{subequations}
  \begin{align}
\langle w_{n-1}(\xi|_{S(F_{b})}),[S(F_{b})]\rangle
&= \chi(\sphere{n-1})\pmod{2}\\
&= 0.
  \end{align}
\end{subequations}
The key poit is that this forces
\begin{equation}
w_{n-1}(\xi|_{S(F_{b})})=0,
\end{equation}
so
\begin{equation}
(j_{b})^{*}(\delta(w_{n-1}(\xi_{0})))=0.
\end{equation}
Then by injectivity of $(j_{b})^{*}$ we find
\begin{equation}
\delta w_{n-1}(\xi_{0})=0,
\end{equation}
so
\begin{equation}
\varphi^{-1}(\delta w_{n-1}(\xi_{0}))=0.
\end{equation}
Hence we have proven our goal, and we can use this procedure to
construct the Stiefel--Whitney class.
\end{node}

\subsection{Properties of the Chern Class}

\begin{proposition}
\begin{enumerate}
\item Naturality: $c(f^{*}\omega)=f^{*}c(\omega)$
\item Product: $c(\omega\oplus\eta)=c(\omega)c(\eta)$
\item $c(\gamma^{1}_{1})\neq0\in H^{2}(\CP^{1};\ZZ)\iso H^{2}(\sphere{2};\ZZ)\iso\ZZ$
\end{enumerate}
\end{proposition}

\begin{proof}[Proof (of (1))]
We want to do induction on the rank of $\omega$.

The inductive hypothesis: assume for all complex rank-$k$ bundles (for
any $k\leq n-1$) the claim holds.

Let $\omega\colon E\xrightarrow{\pi}B$ be a complex rank-$n$ bundle
over $B$. Let $f\colon A\to B$ be continuous. We construct the
pullback bundle $E'$,
\begin{equation}
\vcenter{\xymatrix{
E'\ar[r]\ar[d]^{\pi'} & E\ar[d]^{\pi}\\
A\ar[r]^{f} & B}}
\end{equation}
Then we construct the $E_{0}$, $E'_{0}$ bundles which are related by
the bundle map $f_{0}\colon E'_{0}\to E_{0}$
\begin{equation}
\vcenter{\xymatrix{
E'_{0}\ar[d]\ar@/_1pc/[dd]_{\pi'_{0}}\ar[d]\ar[r]^{f_{0}} & E_{0}\ar[d]\ar@/^1pc/[dd]^{\pi_{0}}\\
E'\ar[r]\ar[d]^{\pi'} & E\ar[d]^{\pi}\\
A\ar[r]^{f} & B}}
\end{equation}
We then construct the bundles $\xi_{0}\to E_{0}$ and $\xi'_{0}\to E'_{0}$
\begin{equation}
\vcenter{\xymatrix{
\xi'_{0}\ar[d] & \xi_{0}\ar[d]\\
E'_{0}\ar[d]\ar@/_1pc/[dd]_{\pi'_{0}}\ar[d]\ar[r]^{f_{0}} & E_{0}\ar[d]\ar@/^1pc/[dd]^{\pi_{0}}\\
E'\ar[r]\ar[d]^{\pi'} & E\ar[d]^{\pi}\\
A\ar[r]^{f} & B}}
\end{equation}
Now we use the assumption that $\omega$ is rank-$n$. Then
$c_{n}(\omega)=e(\omega_{\RR})$, and we have
\begin{subequations}
  \begin{align}
c_{n}(f^{*}\omega) &= e_{n}\bigl((f^{*}\omega)_{\RR}\bigr)\\
&= f^{*}e_{n}(\omega_{\RR})\\
&= f^{*}c_{n}(\omega).
  \end{align}
\end{subequations}
We just need to show that for all $0\leq i\leq n-1$ we have
$c_{i}(f^{*}\omega)=f^{*}c_{i}(\omega)$. Recall
\begin{subequations}
  \begin{equation}
c_{i}(\omega)=(\pi_{0}^{*})^{-1}c_{i}(\omega_{0}),
  \end{equation}
  and
  \begin{align}
c_{i}(f^{*}\omega) &= \bigl((\pi_{0})^{*}\bigr)^{-1}c_{i}(f_{0}^{*}\omega_{0})\\
&=\bigl((\pi_{0})^{*}\bigr)^{-1}f_{0}^{*}c_{i}(\omega_{0})\quad\mbox{by inductive hypothesis},
  \end{align}
  so
  \begin{equation}
f^{*}(\pi_{0}^{*})^{-1}=\bigl((\pi'_{0})^{*}\bigr)^{-1}f_{0}^{*}.
  \end{equation}
\end{subequations}
But by our commutative diagram, specifically relevant part being the
subdiagram:
\begin{equation}
\vcenter{\xymatrix{E'_{0}\ar[d]^{\pi'_{0}}\ar[r]^{f_{0}} & E_{0}\ar[d]^{\pi_{0}}\\
A\ar[r]^{f} & B}}
\end{equation}
we find
\begin{equation}
\pi_{0}\circ f_{0}=f\circ\pi'_{0},
\end{equation}
so
\begin{equation}
(\pi_{0}\circ f_{0})^{*}=(f\circ\pi'_{0})^{*},
\end{equation}
which gives us the claim.
\end{proof}
